\documentclass{beamer}
\usetheme{Berlin}

\title{Fundamentals of Modeling Part 2}
\subtitle{Course: Complex Systems Modeling}
\author{Robert Flassig}
\institute{Brandenburg University of Applied Sciences}
\date{\today}

\begin{document}
\begin{frame}
\titlepage
\end{frame}

\begin{frame}
\frametitle{Outline}
\tableofcontents
\end{frame}

\section{How to create a model}

\begin{frame}{Approaches to Scientific Model Building}
    There are two major families of modeling approaches in science and engineering:
    \begin{itemize}
        \item \textbf{Descriptive Modeling}
        \item \textbf{Rule-Based Modeling}
    \end{itemize}
\end{frame}

\begin{frame}{Descriptive Modeling}
    In descriptive modeling, researchers aim to specify the actual state of a system at a given time. Examples include:
    \begin{itemize}
        \item Taking a picture
        \item Creating a miniature
        \item Writing a biography
    \end{itemize}
    This can also be done using quantitative methods like:
    \begin{itemize}
        \item Regression analysis
        \item Pattern recognition
    \end{itemize}
    Goal: Capture "What the system looks like."
\end{frame}

\begin{frame}{Rule-Based Modeling}
    In rule-based modeling, researchers define dynamical rules to explain system behavior. Examples include:
    \begin{itemize}
        \item Dynamical equations
        \item Theories
        \item First principles
    \end{itemize}
    This approach helps predict future states and can be quantitative or conceptual (e.g., Darwin's evolutionary theory).
    \newline
    Goal: Capture "How the system will behave."
\end{frame}

\begin{frame}{Cycle of Descriptive and Rule-Based Modeling}
    Newton's law of motion is an example of rule-based modeling, making sense of observational information and predicting planetary movement.
    \newline
    Descriptive modeling accumulates descriptions, while rule-based modeling explains them. These two approaches form a cycle in scientific modeling, taking turns in driving discoveries.
\end{frame}

\begin{frame}{Focus on Rule-Based Modeling}
    This course focuses on rule-based modeling, particularly for complex systems. Rule-based models help:
    \begin{itemize}
        \item Understand emergence
        \item Study self-organization
    \end{itemize}
    This is done through computer simulation and mathematical analysis at microscopic scales.
\end{frame}

\begin{frame}{Typical Rule-Based Modeling Cycle}
    The typical steps of rule-based modeling include:
    \begin{enumerate}
        \item Observe the system of interest.
        \item Reflect on possible rules causing the observed characteristics.
        \item Derive predictions and compare with reality.
        \item Modify the rules and repeat the cycle.
    \end{enumerate}
    Termination is based on the modeler's satisfaction or resource limitations.
\end{frame}

\begin{frame}{Challenges in Rule-Based Modeling}
    Among the four steps, reflecting on possible rules is particularly challenging because:
    \begin{itemize}
        \item It relies heavily on the modeler's personal knowledge and experience.
        \item It is deeply connected to how one perceives the world, making it a difficult process to teach or learn.
    \end{itemize}
\end{frame}

\begin{frame}{Examples}
Here are some examples...
\end{frame}

\begin{frame}{Examples of Descriptive vs. Rule-Based Models}
\begin{columns}
\column{0.48\textwidth}
\textbf{Descriptive (What it looks like)}
\begin{itemize}
  \item Regression of temperature vs.\ CO2
  \item Population data fitting
  \item Image-based pattern recognition
  \item Empirical load forecasting
\end{itemize}

\column{0.48\textwidth}
\textbf{Rule-Based (How it behaves)}
\begin{itemize}
  \item Pendulum dynamics
  \item Lotka--Volterra predator--prey
  \item Heat diffusion
  \item Sierpiński carpet (geometric rule)
\end{itemize}
\end{columns}

\vspace{0.5em}
\textit{Upcoming: explore these models interactively in the Jupyter notebook.}\\
\textit{But first, let's dig into key considerations for modelers}
\end{frame}


\section{Modeling Complex Systems: Key Considerations}
\begin{frame}{Modeling Complex Systems: Key Considerations}
    \begin{itemize}
        \item \textbf{Key Questions:} What are the main questions to address in the system?
        \item \textbf{Scale of Description:} At what scale should the system's behaviors be described?
        \begin{itemize}
            \item Define the "microscopic" components and their dynamics.
        \end{itemize}
        \item \textbf{System Structure:} How is the system organized?
        \begin{itemize}
            \item Identify components and their interactions.
        \end{itemize}
    \end{itemize}
\end{frame}

\begin{frame}{Modeling Complex Systems: Dynamics and Changes}
    \begin{itemize}
        \item \textbf{Possible States:} What dynamic states can each component exhibit?
        \item \textbf{Temporal Changes:} How does the system evolve over time?
        \begin{itemize}
            \item Define rules for state changes due to interactions.
            \item Determine how interactions between components change over time.
        \end{itemize}
    \end{itemize}
\end{frame}

\subsection{Exercise}
\begin{frame}{From Examples to Model Refinement}

\textbf{Goal:} Connect real-world phenomena to the process of building and refining models.

\medskip
\textbf{Examples:}
\begin{itemize}
    \item \textbf{Egg boiling time:}  
          Predict yolk softness based on time and temperature.  
          Compare a descriptive fit (empirical cooking chart) to a rule-based model (heat diffusion in a sphere).

    \item \textbf{Traffic flow:}  
          Each driver follows local rules (distance, braking, reaction time) —  
          yet global patterns such as stop-and-go waves emerge.

    \item \textbf{Population dynamics:}  
          Simple birth–death or predator–prey rules lead to rich collective behavior
          (equilibrium, oscillation, collapse).
\end{itemize}

\medskip
\textbf{Exercise – Refining Your Model:}
\begin{itemize}
    \item Start from one of these examples or choose your own system.  
    \item Identify the key components, interactions, and rules.  
    \item Revise your model until it reproduces the main observable behavior.  
    \item Prepare a short reasoning: what makes your model descriptive or rule-based?
\end{itemize}

\end{frame}

\begin{frame}{Example Solutions (1/3): Descriptive vs.\ Rule-Based Models}
\small
\textbf{Egg Boiling Time}
\begin{itemize}
    \item \textbf{Descriptive model:}
          Empirical relation between cooking time $t$, water temperature $T_w$,
          and yolk softness $S$:
          \[
          S = a - b\,e^{-c\,t(T_w - 100)}.
          \]
          Captures observed outcomes without physical reasoning.

    \item \textbf{Rule-based model:}
          Transient heat diffusion inside a spherical egg:
          \[
          \frac{\partial T}{\partial t} = \alpha \nabla^2 T, \qquad
          T(R,t)=T_w,\; T(0,0)=T_0.
          \]
          The egg is “soft” when the yolk center reaches $T(0,t)=T_y$.
\end{itemize}

\vspace{0.5em}
\textit{Interpretation:}  
Descriptive models fit data; rule-based models explain underlying heat transport.
\end{frame}

\begin{frame}{Example Solutions (2/3): Descriptive vs.\ Rule-Based Models}
\tiny
\textbf{Traffic Flow}
\begin{itemize}
    \item \textbf{Descriptive model:}
          Empirical \emph{fundamental diagram} relating density $\rho$ and mean speed $v(\rho)$:
          \[
          q(\rho) = \rho\,v(\rho),
          \quad v(\rho) \approx v_\text{max}\!\left(1 - \frac{\rho}{\rho_\text{max}}\right).
          \]
          Derived from aggregated traffic data.

    \item \textbf{Rule-based model:}
          Microscopic car-following dynamics for vehicle $i$:
          \[
          \frac{dv_i}{dt}
          = a\!\left(1-\frac{v_i}{v_0}\right)
          - b\!\left(\frac{s^*(v_i,\Delta v_i)}{s_i}\right)^2,
          \]
          where each driver adjusts acceleration based on distance and speed difference.
\end{itemize}

\vspace{0.5em}
\textit{Interpretation:}  
Local driving rules create global phenomena such as stop-and-go waves.
\end{frame}

\begin{frame}{Example Solutions (3/3): Descriptive vs.\ Rule-Based Models}
\tiny
\textbf{Population Dynamics}
\begin{itemize}
    \item \textbf{Descriptive model:}
          Logistic growth fitted to measured population data:
          \[
          N(t) = \frac{K}{1 + e^{-r(t - t_0)}}.
          \]
          Describes saturation but not interaction mechanisms.

    \item \textbf{Rule-based model:}
          Lotka--Volterra predator–prey system:
          \[
          \frac{dx}{dt} = \alpha x - \beta x y, \qquad
          \frac{dy}{dt} = \delta x y - \gamma y.
          \]
          Captures feedback between resource and consumer populations.
\end{itemize}

\vspace{1em}
\textbf{Takeaway}
\begin{itemize}
    \item Descriptive $\Rightarrow$ summarizes observed behavior (\emph{what is}).
    \item Rule-based $\Rightarrow$ encodes mechanisms (\emph{how it changes}).
    \item Model refinement connects both: fit parameters $\leftrightarrow$ test rules.
\end{itemize}

\end{frame}



\end{document}